% 第八批：原 PDF 物理页420–423中部。
\chapter{结语}\label{chap:conclusion}
\SourcePage{420}
现在已来到本书的结尾。一路上见识了一些精彩的数学，也希望已对镜像对称现象有所了解。在这最后一章，将回顾已完成的工作及尚待完成的工作，还会简要讨论前几章未涉及的镜像对称的众多方面中的若干主题。

\section{回顾与未解问题}\label{sec:reflections-open-problems}
研究镜像对称的过程中，用到了数量惊人的代数几何知识：讨论过环面簇、Hodge 结构变分、混合 Hodge 结构、具有正则奇点的微分方程、超几何方程、Frobenius 方法、Kähler 锥、种种模空间、辛几何、代数叠、计数几何、Clemens 猜想、等变上同调和局部化，以上还只是其中一部分。镜像对称把代数几何的如此多部分汇聚起来，表明我们共同从事的这项事业，其内在联系比我们所认识的更为深刻。

不过，镜像对称不仅是集中运用已有数学。在本书中，还见证了新对象进入代数几何研究，包括 Kähler 模、稳定映射、Gromov--Witten 不变量、量子上同调、A 模型 Hodge 结构变分、量子微分方程、Euler 数据以及上同调值形式函数。把它们与物理学中一些尚未有严密数学定义的思想结合起来，就能窥见二十一世纪代数几何可能带给我们的内容。

还要强调，镜像对称这一领域比本书所呈现的丰富得多。除了几乎未触及的庞大物理文献，还有许多数学方面完全没有涉及。第 \hyperref[sec:other-aspects-mirror-symmetry]{\ref*{sec:other-aspects-mirror-symmetry}} 节将列出其中一部分。不过，即使只限于前几章呈现的镜像对称观点，仍有很多工作尚待完成，下面就来说明。

\subsection{什么是镜像对称？}\label{subsec:what-is-mirror-symmetry}
基本问题在于，尚未提出一个完全令人满意的镜像定理版本，甚至尚未提出这样的镜像猜想。回顾迄今的工作，可辨认出三种不同的镜像对称形式：
\begin{itemize}
\item \hyperref[sec:mirror-symmetry-revisited]{第~\ref*{sec:mirror-symmetry-revisited} 节}探讨了 \SourceCite{Morrison7} 的建议：数学版本的镜像对称，应表述为一个光滑
\SourcePage{421}
Calabi--Yau 流形的 A 模型 Hodge 结构变分与其镜像的几何（即 B 模型）Hodge 结构变分之间的同构。
\item \hyperref[sec:quintic-mirror-theorem]{第~\ref*{sec:quintic-mirror-theorem} 节}研究了 \SourceCite{LLY} 的工作，它把镜像定理表述为 Euler 数据 $\widehat P$ 与 $\widehat Q$ 在镜像变换后的等价。这意味着某些超几何的上同调值函数之间存在关系，从中可计算五次三维簇的 Gromov--Witten 势。
\item \hyperref[sec:givental-mirror-approach]{第~\ref*{sec:givental-mirror-approach} 节}讨论了 \SourceCite{Givental2,Givental4} 的定理，其中镜像定理表述为上同调值函数 $I_{\mathcal V}$、$J_{\mathcal V}$ 在镜像变换后相等。所得等式给出了五次三维簇的 Gromov--Witten 势的一个等价描述。
\end{itemize}
\hyperref[sec:quintic-other-methods]{第~\ref*{sec:quintic-other-methods} 节}预览过这三种方法，它们在 $\mathbb P^4$ 中五次三维簇的情形下巧妙地汇合。不过，这些方法也有明显差别，例如：
\begin{itemize}
\item \hyperref[sec:mirror-symmetry-revisited]{第~\ref*{sec:mirror-symmetry-revisited} 节}的 Hodge 理论方法只对 Calabi--Yau 流形有意义。
\item \SourceCite{LLY} 的方法适用于多种情形，包括五次三维簇，以及\hyperref[subsec:multiple-cover-localization]{第~\ref*{subsec:multiple-cover-localization} 小节}研究的多重覆盖公式。但在本书写作时，这些结果的环境空间仍限于 $\mathbb P^n$。\footnote{我们刚获悉（1998 年 11 月），\SourceCite{LLY} 的作者即将宣布该文内容的一个环面版本。预印本发表后，应能得到环面镜像定理的详细证明。}
\item \SourceCite{Givental2,Givental4} 的定理适用于光滑环面簇 $X_\Sigma$ 中的 nef 完全交 $X$。$X=X_\Sigma$ 时也有非平凡结果，但这些方法例如不适用于多重覆盖公式。
\end{itemize}
后两种方法令人兴奋，因为终于能够实际证明镜像对称中的一些漂亮公式。回顾\hyperref[chap:mirror-theorem]{第十一章}镜像定理的证明中如何使用等变上同调及局部化，就能看出其局限。这些方法需要群作用，这解释了环面簇为何起着如此突出的作用。不过，结果不仅适用于 Calabi--Yau 情形。因此，\SourceCite{LLY} 与 \SourceCite{Givental2} 的方法可能属于一个更广泛的“环面”版本的镜像对称，Calabi--Yau 环面完全交则是其特例。

另一方面，\hyperref[sec:mirror-symmetry-revisited]{第~\ref*{sec:mirror-symmetry-revisited} 节}的 Hodge 理论版本与环面簇无关。任意射影 Calabi--Yau 流形 $V$ 都有 A 模型 Hodge 结构变分，可以问：它是否与某个 Calabi--Yau 流形的几何 Hodge 结构变分相同？若是，后者便是镜像 $V^\circ$。此外，如例题~\ref{ex:octic-toric-mirror-theorem} 后的评论所示，Hodge 理论版本的镜像对称与 Givental 版本的镜像定理之间的关系，仍有待弄清。最后，如第 \hyperref[subsec:homological-categorical-mirror]{\ref*{subsec:homological-categorical-mirror}} 小节将展示的，有些镜像对称版本涉及的内容远远超出前几章。所有这些都表明，距离完全理解镜像对称的数学，我们仍很遥远。

最后，提醒读者：按照物理学，镜像对称是两个物理理论的等价，而这些理论目前尚未得到严密的数学定义。
\SourcePage{422}
随着这些理论被置于更牢固的数学基础之上，将能更深入地理解镜像对称的数学。这可能填补刚才提到的一些空缺，也可能把我们引向完全出乎意料的方向。

\subsection{尚未解决的猜想}\label{subsec:open-conjectures}
如刚解释的，镜像对称中主要的未解问题是提出并证明一般的镜像定理。与此有关，现有镜像定理版本也有一些有趣问题。例如，应更好地理解\hyperref[subsec:givental-toric-mirror]{第~\ref*{subsec:givental-toric-mirror} 小节} Givental 的环面镜像定理版本，与\hyperref[subsec:toric-mirror-conjecture]{第~\ref*{subsec:toric-mirror-conjecture} 小节}的环面镜像猜想之间的联系。例题~\ref{ex:octic-toric-mirror-theorem} 结合此结果与式~\eqref{eq:octic-virtual-class-pushforward-assumption}，导出了式~\eqref{eq:octic-toric-mirror-potential} 和~\eqref{eq:octic-yukawa-period-second-derivative}。我们很希望知道，这些公式可以推广到何种程度。

不过，镜像定理提出的问题并非全部。读者还应记得：为了陈述本书所考虑的镜像对称版本，在建立所需工具的途中，也遇到了一些有趣的未解问题。现在回顾其中一些猜想，希望鼓励人们研究它们。

先列出\hyperref[chap:hodge-yukawa]{第五章}和\hyperref[chap:moduli]{第六章}讨论的、关于复模与 Kähler 模的四个猜想：
\begin{itemize}
\item \hyperref[subsec:higher-dimensional-mum]{第~\ref*{subsec:higher-dimensional-mum} 小节}的\emph{整性猜想}描述了最大幺幂边界点处，权滤过的整数结构应如何与单值性相互作用。
\item \hyperref[subsec:toric-complex-moduli]{第~\ref*{subsec:toric-complex-moduli} 小节}在 Calabi--Yau 环面超曲面的简化模空间上构造了某些特殊边界点。猜想~\ref{conj:toric-mum-boundary} 断言，它们是最大幺幂边界点。
\item \hyperref[subsec:large-radius-limits]{第~\ref*{subsec:large-radius-limits} 小节}的\emph{锥猜想}断言，维数大于二的 Calabi--Yau 流形 $V$ 的 Kähler 锥，模去 $\operatorname{Aut}(V)$ 的作用后应为有理多面锥。
\item \hyperref[subsec:enlarged-kahler-moduli]{第~\ref*{subsec:enlarged-kahler-moduli} 小节}的猜想~\ref{conj:toric-kahler-cone} 描述 Calabi--Yau 环面超曲面的 Kähler 锥的环面部分。
\end{itemize}
这些猜想自身都很有趣，同时也是镜像对称的重要组成部分。

除了这些模空间猜想，还有许多与镜像对称有关的未解问题值得代数几何学家思考。以下是我们觉得特别有趣的一些：
\begin{itemize}
\item \hyperref[subsec:toric-gkz]{第~\ref*{subsec:toric-gkz} 小节}研究 GKZ 方程时，看到这些方程给出了部分 Picard--Fuchs 方程，却并非全部。应当能够描述缺少的 Picard--Fuchs 方程。此外，上同调值形式函数 $I_{\mathcal V}$ 应满足镜像的 Picard--Fuchs 方程。
\item \hyperref[subsec:normalized-yukawa-connection]{第~\ref*{subsec:normalized-yukawa-connection} 小节}讨论了最大幺幂边界点处 Calabi--Yau 三维簇的归一化 Yukawa 耦合与势函数。当三维簇为环面完全交时，应能像例题~\ref{ex:octic-toric-mirror-theorem} 那样，用适当的 $I_{\mathcal V}$ 计算 Yukawa 耦合和势函数。
\item 关于虚拟基本类 $[\overline M_{g,n}(X,\beta)]^{\mathrm{virt}}$，有一些问题值得进一步研究。\hyperref[subsec:gw-equivalence]{第~\ref*{subsec:gw-equivalence} 小节}讨论了 Gromov--Witten 不变量与虚拟基本类的各种定义之间的等价性。另一个问题是，当 $X$ 为 $Y$ 中的完全交时，$[\overline M_{g,n}(X,\beta)]^{\mathrm{virt}}$ 与
\SourcePage{423}
$[\overline M_{g,n}(Y,\beta)]^{\mathrm{virt}}$ 有何联系。式~\eqref{eq:octic-virtual-class-pushforward-assumption} 表明，这对理解 Givental 函数 $J_X$ 与 $J_{\mathcal V}$ 的关系很重要。
\item \IndexLocation{indexloc:1131:423}量子上同调的定义涉及无穷和，\hyperref[subsec:small-quantum-definition]{第~\ref*{subsec:small-quantum-definition} 小节}猜测，此和在大半径极限点的某个邻域中收敛。注意，只有此猜想成立，A 模型 Hodge 结构变分才有意义。
\end{itemize}
如前，即使不考虑与镜像对称的联系，我们也认为这些问题很有趣。

本书中一些最重要的未解问题涉及计数几何。起初，\SourceCite{CdGP} 对五次三维簇的预言近乎神奇：物理学似乎毫不费力地解决了远超代数几何标准方法能力的计数问题。此后，镜像对称已成为计算 Gromov--Witten 不变量的有力工具，但理解一个 Gromov--Witten 不变量的计数含义，仍是非平凡的问题。即使对于五次三维簇 $V$，瞬子数 $n_{10}$ 与 $V$ 上十次有理曲线的关系，也可能很复杂，因为十次情形的 Clemens 猜想尚未获证。即使获证，$n_{10}$ 也不是这类曲线的数目，如\hyperref[subsec:instanton-numbers]{第~\ref*{subsec:instanton-numbers} 小节}所解释。另一个很好的问题是猜想~\ref{conj:symplectic-instantons}，它试图直接定义瞬子数，而不是像本书那样，间接地通过 Gromov--Witten 不变量 $\langle I_{0,0,\beta}\rangle$ 定义它们。不过，瞬子数的问题只是全貌的一部分。如\hyperref[subsec:instanton-numbers]{第~\ref*{subsec:instanton-numbers} 小节}和\hyperref[sec:computing-gw-ii]{第~\ref*{sec:computing-gw-ii} 节}的例题及参考文献所示，镜像对称使计数几何论文数量激增，其中许多与 Calabi--Yau 流形或 Gromov--Witten 不变量无关。

如前言所说，本书应视为对镜像对称这一仍在进行中的研究的报告。这个领域的某些部分已有牢固基础，其他部分仍在迅速变化。镜像对称的权威性著作尚有待写出。不过，希望本书所呈现的虽有限的镜像对称观点，已经传达了代数几何这一惊人领域的一些奇妙与激动人心之处。

% 第九批从原 PDF 423 中部接续；该页上部锚点已在第八批。
\hypertarget{source:423-batch9}{}
\section{镜像对称的其他方面}\label{sec:other-aspects-mirror-symmetry}
主要目标现已完成，但还有许多镜像对称的方面，本书未曾涉及。在本书写作时，镜像对称仍是活跃的研究领域，许多新思想正不断发展，并由此带来数学与物理学中的其他新进展。这里甚至不打算给出完整清单，只举几个例子。下面讨论的主题，有些性质纯属数学，有些则与物理学联系更密切，但它们都包含值得关注的数学内容。

\subsection{K3 曲面上的有理曲线}\label{subsec:rational-curves-k3}
固定一个一般的 K3 曲面，其中有除子类 $C$ 满足 $C^2=2g-2$。于是，$|C|$ 是一个 $g$ 维线性系，其中的曲线具有算术亏格 $g$。令 $n(g)$ 为这个线性系中有理曲线的数目。
\SourcePage{424}
Yau 与 Zaslow \SourceCite{YZ} 提出了如下惊人的公式猜想：
\[
1+\sum_{g=1}^{\infty}n(g)q^g=q/\Delta(q),
\]
其中 $\Delta(q)$ 是通常的权为 $12$ 的模形式。最早的数学阐述见 \SourceCite{Beauville2}，其中也证明了若干中间步骤。基本想法是：对这个线性系中出现的每条有理曲线，赋予一个重数，它等于该曲线的紧化 Jacobian 的 Euler 示性数（同一个 Euler 示性数也出现在 \SourceCite{YZ} 中）。这个重数可以与稳定映射联系起来 \SourceCite{FGvS}。此外，\SourceCite{FGvS} 还证明，此重数总为正数。因此，Yau--Zaslow 公式有明确的计数含义，即按适当重数计数有理曲线。

\SourceCite{Gottsche} 把 Yau--Zaslow 公式推广为一个关于代数曲面上结点曲线的猜想，提出一个生成函数，用以计数曲线的 $n$ 维族中具有 $n$ 个结点的曲线。当曲面是 K3 曲面或 Abel 曲面时，猜想中的生成函数可用拟模形式表示。对 K3 曲面，这个猜想已在 \SourceCite{BL} 中得到证明；下面稍加说明。考虑满足 $C^2=2g+2n-2$ 的曲线线性系 $|C|$。它的维数为 $g+n$，其中的曲线具有算术亏格 $g+n$。预期 $|C|$ 中具有 $n$ 个结点的曲线构成一个 $g$ 维子族。令 $N_g(n)$ 表示这个线性系中经过 $g$ 个一般点、几何亏格为 $g$ 且有 $n$ 个结点的曲线的数目。\SourceCite{BL} 证明了
\[
\sum_{n=0}^{\infty}N_g(n)q^n
=\left(\frac{d}{dq}G_2(q)\right)^g\frac{q}{\Delta(q)},
\]
其中 $G_2(q)$ 是 Eisenstein 级数
\[
G_2(q)=-\frac1{24}+\sum_{k=1}^{\infty}\sigma(k)q^k,
\]
$\sigma(k)$ 是算术函数 $\sum_{d\mid k}d$。注意，$G_2$ 及其导数都是拟模函数。
\index{K3曲面@K3 曲面!有理曲线|hyperpage}\index{Yau--Zaslow公式@Yau--Zaslow 公式|hyperpage}\index{拟模形式|hyperpage}

\subsection{K3 曲面的镜像对称与 Arnold 奇异对偶性}\label{subsec:k3-arnold-strange-duality}
二十世纪七十年代末，不同研究者指出，K3 曲面的一种对偶性可以解释例外单模临界点的 Arnold 奇异对偶性。近年人们认识到，K3 曲面的这种对偶性属于二维版本的镜像对称 \SourceCite{Dolgachev3,Kobayashi,Ebeling,GN1,GN3}。

K3 曲面的镜像对称有不同的特点。对维数大于二的 Calabi--Yau 流形 $V$，Kähler 锥在 $H^2(V,\mathbb R)$ 中是开集；对 K3 曲面 $S$，情形便不再如此，而且 Kähler 锥与格 $H^2(S,\mathbb Z)$ 之间的关系相当微妙。因此，格论在 K3 曲面的镜像对称中起重要作用，并不令人意外。\hyperref[sec:voisin-borcea]{第~\ref*{sec:voisin-borcea} 节}给出的 Voisin--Borcea 构造，已让我们见到这一点的一些迹象。

也可以在 K3 曲面 $S$ 上研究弦理论，\SourceCite{AM3} 对此作了说明。此时仍有 SCFT 模空间，但由于 $h^{2,0}(S)\ne0$，复结构与 Kähler 结构的形变不再解耦，因此不再得到
\SourcePage{425}
式~\eqref{eq:scft-foliations} 所示的局部乘积结构。这一主题内容广泛，无法在几个段落中概述。关于更多细节及相关文献，请参阅 \SourceCite{AM3,Aspinwall}。
\index{Arnold奇异对偶性@Arnold 奇异对偶性|hyperpage}\index{K3曲面@K3 曲面!镜像对称|hyperpage}

\subsection{McKay 对应}\label{subsec:mckay-correspondence}
若 $G\subset\operatorname{SL}(n,\mathbb C)$ 是有限群，且 $\mathbb C^n/G$ 有一个满足 $K_Y=0$ 的消解 $Y$，则 \emph{McKay 对应} \SourceCite{Reid6} 断言，$G$ 的不可约表示应给出 $H^*(Y,\mathbb Z)$ 的一组基，使 $G$ 的特征标表对应于 $Y$ 上的 Poincaré 对偶性。此猜想对物理学中的 Euler 数 \SourceCite{Vafa1,Roan2}，以及 Batyrev 与 Dais 的弦论 Hodge 数 \SourceCite{BD}，都有相应推论。
\index{McKay对应@McKay 对应|hyperpage}

\subsection{超 Kähler 流形}\label{subsec:hyperkahler-manifolds}
若紧致复流形 $M$ 有一个闭的、非退化的全纯二形式，就称 $M$ 是\emph{全纯辛}的。由 Calabi--Yau 定理，$M$ 有一个度量及一个关于 Riemann 联络平行的四元数作用，因此 $M$ 是超 Kähler 流形。\SourceCite{Verbitsky} 中，Verbitsky 证明，超 Kähler 流形是它自身的镜像。
\index{全纯辛流形|hyperpage}\index{超Kähler流形@超 Kähler 流形|hyperpage}

\subsection{D 膜与 T 对偶}\label{subsec:d-branes-t-duality}
本书将镜像对称理解为与镜像对 $(V,\omega)$ 和 $(V^\circ,\omega^\circ)$ 相联系的两个 SCFT 相等（见\hyperref[chap:introduction]{第一章}的讨论）。SCFT 中的态称为\emph{微扰态}，其中包括与 $V$ 和 $V^\circ$ 上的上同调类对应的态。近年人们认识到，这个理论在某种意义上可以扩充，以包含\emph{非微扰态}，包括与 \emph{D 膜}相联系的态。D 膜本质上是一个 $p$ 维物体在时空中传播时扫出的\emph{世界体}。为明确维数，有时将它们称为 D$p$ 膜；由于多了一个时间维，它们的世界体是 $p+1$ 维的。例如，D1 膜的样子与弦的世界面完全相同（如\hyperref[chap:introduction]{第一章}所绘）。态本身来自边界落在 D 膜上的开弦。D 膜的综述见 \SourceCite{Polchinski1}。人们猜测，与 $(V,\omega)$ 和 $(V^\circ,\omega^\circ)$ 相联系的物理理论，\emph{连同非微扰态在内}都相等，从而使镜像对称成为更强的断言。

\SourceCite{SYZ} 中，Strominger、Yau 与 Zaslow 提出，镜像对称就是物理学家所说的 \emph{T 对偶}。下面简述这个想法。注意，$V^\circ$ 上点的模空间与 $V^\circ$ 本身同构，而它又自然等同于 $V^\circ$ 上 D0 膜模空间的一部分。为理解这一点，回顾\hyperref[chap:introduction]{第一章}：弦理论是十维理论，其中 $V^\circ$ 贡献六个维度；时间维是余下四个维度之一。对任一点 $p\in V^\circ$，当 $p$ 随时间传播时，便扫出一条 D0 膜的世界线。镜像对称预言，它应对应于 $V$ 上 D3 膜模空间的一部分。进一步分析表明，Calabi--Yau 流形 $V$ 的镜像 $V^\circ$，应是 $V$ 上某些特殊 Lagrange 环面的模空间，经过适当复化与紧化得到的空间。如 \SourceCite{Morrison8} 所解释，T 对偶或许最终可以这样给出一个纯数学的镜像对称构造。在这个描述中，$V$ 与 $V^\circ$ 都映到同一个实三维底空间，一般纤维是实三维环面。人们猜测，$V$ 与 $V^\circ$ 的这两个纤维化中的对应纤维互为对偶环面，“T 对偶”之名也由此而来。这可能导出镜像流形的拓扑构造。上述断言的某些方面，已在 Voisin--Borcea 三维簇 \SourceCite{GW} 和 K3 曲面
\SourcePage{426}
\SourceCite{GW,Morrison8} 的情形下得到检验。对椭圆 K3 曲面，T 对偶与 Fourier--Mukai 变换有关 \SourceCite{BBRP}。\SourceCite{Gross3} 中还提出了若干值得关注的拓扑推论猜想。
\index{D膜@D 膜|hyperpage}\index{T对偶@T 对偶|hyperpage}\index{世界体|hyperpage}\index{特殊Lagrange环面@特殊 Lagrange 环面|hyperpage}

\subsection{同调镜像对称或范畴镜像对称}\label{subsec:homological-categorical-mirror}
1994 年，Kontsevich \SourceCite{Kontsevich3} 提出了镜像对称的“同调”版本，将镜像猜想表述为范畴等价。此时，$V^\circ$ 是 $V$ 的镜像，意味着 $V^\circ$ 上凝聚层的导出范畴，应等价于 $V$ 上带酉局部系统的 Lagrange 子流形的导出范畴。由 Lagrange 子流形构造一个范畴的基本想法来自 \SourceCite{Fukaya}。\SourceCite{PZ} 建议，或许应限制到 $V$ 的特殊 Lagrange 子流形。注意，Lagrange 子流形使用 $V$ 的 Kähler 结构（A 模型），凝聚层则使用 $V^\circ$ 的复结构（B 模型）。

\SourceCite{Kontsevich3} 解释，与 $V$ 相联系的 Lagrange 范畴的某些形变，应当使我们能够将 $V$ 的量子上同调构造为通常杯积的形变。这里的形变自然地在整个 $H^*(V)$ 上进行，而非像\hyperref[sec:a-model-vhs]{第~\ref*{sec:a-model-vhs} 节} A 模型联络的定义那样，只在 $H^{1,1}(V)$ 上进行。由于它应与 $V^\circ$ 上凝聚层的形变理论有关，这暗示着一种超出通常复结构形变的丰富结构。下面的公式可以让我们窥见其中可能涉及的内容：
\begin{equation}\label{eq:hochschild-polyvector-cohomology}
HH^n(V^\circ)=\bigoplus_{p+q=n}H^q(V^\circ,\wedge^p\Theta_{V^\circ}),
\end{equation}
其中 $HH$ 表示 Hochschild 上同调。这个上同调群应与 $V^\circ$ 的复模空间的某个加厚空间的切空间等同。注意，式~\eqref{eq:hochschild-polyvector-cohomology} 的各直和项，恰好是式~\eqref{eq:eigenspaces} 第一行所写 $N=2$ 超共形代数中算子 $(Q,\overline Q)$ 的 $(p,q)$ 特征子空间。以式~\eqref{eq:hochschild-polyvector-cohomology} 为切空间来加厚模空间的想法，最早出现在 \SourceCite{Witten4}。这些空间能够在数学中自然出现，是一件很美妙的事。如 \SourceCite{Kontsevich3} 所解释，$HH^2(V^\circ)$ 包含 $V^\circ$ 的复结构的通常一阶形变、代数层 $\mathcal O_{V^\circ}$ 的某些非交换形变，以及从形变观点看仍颇为神秘的一些项。由于维数至少为三的 Calabi--Yau 流形满足 $H^2(V^\circ,\mathcal O_{V^\circ})=0$ 和 $H^0(V^\circ,\wedge^2\Theta_{V^\circ})=0$，最后这两项在此时不成问题；不过，$n\ne2$ 时 $HH^n(V^\circ)$ 的含义仍有些神秘。此外，\SourceCite{Kontsevich3} 指出，猜想 $V$ 的 Lagrange 范畴具有 $A_\infty$ 结构。（不过，$A_\infty$ 范畴并非通常意义下的范畴，其复合不再严格满足结合律，因而须谨慎处理。）这再次暗示，在镜像一侧，$V^\circ$ 上凝聚层的范畴应有更丰富的结构。这些都表明，我们尚未发现镜像对称的全部内容。

\SourceCite{PZ} 最近给出了椭圆曲线情形下这种范畴等价的一个例子。该文还指出，T 对偶可能是同调镜像对称的推论。在这个背景下，给定一点 $p\in V^\circ$，子簇 $\{p\}\subset V^\circ$ 的凝聚层 $\mathcal O_{\{p\}}$，应对应于 $V$ 的某个特殊 Lagrange 子流形。这是
\SourcePage{427}
\hyperref[subsec:d-branes-t-duality]{第~\ref*{subsec:d-branes-t-duality} 小节}所述思想的一个版本：镜像 $V^\circ$ 是 $V$ 中某些实 Lagrange 环面的模空间。
\index{同调镜像对称|hyperpage}\index{范畴镜像对称|hyperpage}\index{Hochschild上同调@Hochschild 上同调|hyperpage}\index{A无穷范畴@$A_\infty$ 范畴|hyperpage}

\subsection{K3 纤维化与椭圆纤维化的 Calabi--Yau 流形}\label{subsec:k3-elliptic-fibered-cy}
具有 K3 纤维化或椭圆纤维化的 Calabi--Yau 流形，是大量数学与物理论文的研究对象。从数学上说，对这些纤维化感兴趣相当自然。在研究 Calabi--Yau 流形上的线性系时，K3 纤维化与椭圆纤维化起特殊作用，因为它们自然地出现在 Calabi--Yau 流形的对数极小模型纲领中。在通常的极小模型纲领中，人们试图把簇理解为纤维是极小模型的纤维化。对 Calabi--Yau 簇的对数极小模型，相应情形是纤维为较低维的 Calabi--Yau 簇。\SourceCite{Oguiso} 已开始对这些纤维化进行分类。

人们也研究了 K3 纤维化的镜像对称，猜测纤维的镜像应影响总空间的镜像。\SourceCite{GN2} 讨论了这个想法，\SourceCite{AKMS,HLY3} 中可找到这类空间的例子。

K3 纤维化与椭圆纤维化也出现在物理学中，但原因不同。推动研究 K3 纤维化的主要动力是\emph{对偶性}。更具体地说，人们相信，与 Calabi--Yau 三维簇相联系的 \emph{II 型}弦理论，对偶于与 K3 曲面和椭圆曲线的乘积相联系的杂化弦理论。这意味着，各理论的参数取某些值时，两个理论成为同一个理论。K3 曲面上的向量丛，是定义这类特定杂化弦理论所需数据的一部分。对这种对偶性，已有大量数学检验。

椭圆纤维化的 Calabi--Yau 三维簇，在物理学中是可以用来构造 \emph{F 理论}模型的空间 \SourceCite{Vafa3}。由于丛在 F 理论中也起关键作用，这促使人们进一步研究具有椭圆纤维化的 Calabi--Yau 流形上的主丛 \SourceCite{FMW1}。关于这些丛的数学研究，见 \SourceCite{Donagi,FMW2}。与此密切相关的是 F 理论和杂化弦理论之间的对偶性，它出现在由同时具有 K3 纤维化和椭圆纤维化的 Calabi--Yau 三维簇得到的理论中。这种对偶性可以视为低一维情形下类似结果的推广。与椭圆纤维化的 K3 曲面相联系的 F 理论，和椭圆曲线上的杂化弦理论之间的对偶性，有一个优美的数学支持：椭圆曲线上丛的模，可与某些类型的椭圆纤维化 K3 曲面的 Weierstrass 模型的模等同。
\index{K3纤维化@K3 纤维化|hyperpage}\index{椭圆纤维化|hyperpage}\index{杂化弦理论|hyperpage}\index{F理论@F 理论|hyperpage}

\subsection{锥形奇点跃迁}\label{subsec:conifold-transitions-conclusion}
当 Calabi--Yau 三维簇出现结点奇点，并有一个仍为 Calabi--Yau 的消解时，就发生了\emph{锥形奇点跃迁}。由杂化弦得到的 SCFT 会在这一跃迁过程中变得奇异。不过，\SourceCite{GMS} 表明，\emph{II 型}弦理论可以避免锥形奇点跃迁处的奇异性。粗略地说，与原流形中的一个消失圈（把它看作一个三维膜\footnote{\hyperref[subsec:d-branes-t-duality]{第~\ref*{subsec:d-branes-t-duality} 小节}所考虑的 D3 膜，是三维膜的一种特殊类型。直到 \SourceCite{GMS} 写成之后，D 膜才得到普遍接受。现在人们一致认为，\SourceCite{GMS} 中的三维膜实际上就是 D3 膜。}）相联系的带电黑洞，在 Calabi--Yau 簇变得奇异时质量变为零；这一奇异性补偿了前面所说理论中的奇异性。随后，这些黑洞在消解上成为基本粒子。若假定 Reid 幻想 \SourceCite{Reid5} 的一个 Kähler 版本成立，
\SourcePage{428}
便可以连接任意两个 Calabi--Yau 三维簇所产生的物理理论。这已在作为加权射影空间超曲面的 Calabi--Yau 三维簇情形下得到验证 \SourceCite{ACJM,CGGK}，尽管其奇点可能不是结点。\hyperref[subsec:kahler-moduli-boundary]{第~\ref*{subsec:kahler-moduli-boundary} 小节}讨论过锥形奇点跃迁与镜像对称有关的某些方面。

\SourceCite{Morrison9} 猜测，锥形奇点跃迁的镜像也是锥形奇点跃迁；更一般地，极端跃迁的镜像也是极端跃迁。\SourceCite{BCFKvS1,BCFKvS2} 已在一些例子中检验了这一点。

\subsection{旗流形中的 Calabi--Yau 完全交}\label{subsec:flag-cy-complete-intersections}
这里的关键观察来自 \SourceCite{Sturmfels2}：Grassmann 簇 $G(k,n)$ 可以平坦形变到一个射影环面簇 $P(k,n)$。此外，根据 \SourceCite{BCFKvS2}，$P(k,n)$ 是 Fano 簇，即它来自一个自反多胞形。利用这一点，\SourceCite{BCFKvS1} 的作者将\hyperref[chap:mirror-constructions]{第四章}的 Batyrev--Borisov 镜像构造，与\hyperref[subsec:kahler-moduli-boundary]{第~\ref*{subsec:kahler-moduli-boundary} 小节}及\hyperref[subsec:conifold-transitions-conclusion]{第~\ref*{subsec:conifold-transitions-conclusion} 小节}讨论的极端跃迁结合起来，给出 $G(k,n)$ 中 Calabi--Yau 完全交 $V$ 的镜像构造。

基本想法如下。设 $V\subset G(k,n)$ 是次数为 $d_1,\ldots,d_r$ 的超曲面的完全交。将 $G(k,n)$ 退化到 $P(k,n)$ 时，$V$ 退化到环面完全交 $\overline V\subset P(k,n)$。利用光滑的保典范消解 $\widehat P(k,n)\to P(k,n)$，可得到一个消解 $V'\to\overline V$。若将它理解为先从 $V'$ 经收缩到 $\overline V$，再从 $\overline V$ 经光滑化到 $V$，就得到了\hyperref[subsec:kahler-moduli-boundary]{第~\ref*{subsec:kahler-moduli-boundary} 小节}讨论的一种极端收缩（此时，由于 $\overline V$ 的奇点尤其温和，这一跃迁是锥形奇点跃迁）。如该小节所述，镜像对称暗示应当有一个“镜像”极端跃迁：$V'$ 的镜像 $V'^\circ$ 退化到 $\overline V^\circ$，然后消解为 $V$ 的镜像 $V^\circ$。

在此情形中，$V'\subset\widehat P(k,n)$ 的镜像 $V'^\circ$，可由\hyperref[sec:toric-complete-intersections]{第~\ref*{sec:toric-complete-intersections} 节}所述的 Batyrev--Borisov 镜像构造得到。通过求出相应的退化与光滑化，\SourceCite{BCFKvS1} 能够明确描述 Grassmann 簇中 Calabi--Yau \IndexLocation{indexloc:0802:428}完全交的镜像构造。

为利用这个构造计算 Gromov--Witten 不变量，\SourceCite{BCFKvS1} 引入了\emph{阶乘技巧}；该文将它描述为\hyperref[subsec:quantum-hyperplane-section-principle]{第~\ref*{subsec:quantum-hyperplane-section-principle} 小节}的\IndexLocation{indexloc:1159:428}量子超平面截面原理的一个朴素版本。\SourceCite{BCFKvS1} 使用的另一工具，是量子上同调与量子微分方程之间的关系，我们已在\hyperref[sec:quantum-cohomology-relations]{第~\ref*{sec:quantum-cohomology-relations} 节}讨论过。由于 Grassmann 簇的量子上同调已知，就能获得关于量子微分方程的大量信息（\SourceCite{BCFKvS1} 将其称为\emph{量子 $\mathcal D$ 模}）。再利用\hyperref[subsec:givental-toric-mirror]{第~\ref*{subsec:givental-toric-mirror} 小节}的\IndexLocation{indexloc:0862:428}环面镜像定理和 \SourceCite{Kim2} 中给出的推广，便可计算 Grassmann 簇中许多三维 Calabi--Yau 完全交的瞬子数。

上述工作都可在更一般的情形下进行。令 $F(n_1,\ldots,n_\ell,n)$ 表示由所有满足
\[
0\subset W_1\subset\cdots\subset W_\ell\subset\mathbb C^n,
\qquad \dim W_i=n_i
\]
的旗组成的部分旗流形。\SourceCite{GL} 证明，$F(n_1,\ldots,n_\ell,n)$ 退化到一个环面簇 $P(n_1,\ldots,n_\ell,n)$，推广了上面提到的 Sturmfels 的结果。由于 \SourceCite{BCFKvS2} 表明 $P(n_1,\ldots,n_\ell,n)$ 是 Fano 簇，以上构造便可应用。还应提及 \SourceCite{Batyrev5}，它回顾了刚才描述的构造，并考察一般情形下 Fano 簇退化到 Fano 环面簇时会发生什么。
\index{旗流形!Calabi--Yau完全交@Calabi--Yau 完全交|hyperpage}\index{阶乘技巧|hyperpage}\index{量子D模@量子 $\mathcal D$ 模|hyperpage}

\SourcePage{429}
\subsection{物理学中的其他主题}\label{subsec:other-physics-topics}
镜像对称的物理学近来有一长串新进展，是前几章没有讨论的。其中一些已在本节提到，但还有许多。例如，物理学论证使人们预期：当 Calabi--Yau 流形中的某些曲线被收缩时，可将超对称场论作为弦理论的一个极限得到。因此，Calabi--Yau 几何提供了分析这些超对称理论的工具，尤其是分析其规范理论方面的工具。这些场论有可利用镜像对称求解的“周期”，从而例如可以导出 Seiberg--Witten 曲线及其微分，而无需假定任何对偶猜想。这方面的镜像对称应用见 \SourceCite{KKV,KMV}；利用 Calabi--Yau 几何构造量子场论的其他情形见 \SourceCite{KV,BJPSV,BV}。这样的理论称为\emph{几何工程理论}，其入门介绍见 \SourceCite{Mayr}。在数学上，几何工程理论明确预言了\hyperref[subsec:kahler-moduli-boundary]{第~\ref*{subsec:kahler-moduli-boundary} 小节}所讨论的极端跃迁中光滑化分支的维数。最简单的例子是锥形奇点跃迁，此时 \SourceCite{GMS} 中来自物理学的预言与数学相符；这一情形的光滑化维数，可根据 \SourceCite{Clemens} 作数学计算。在本书写作时，物理学所给一般预言的正确性尚无数学证明。\SourceCite{KMP} 已在一些例子中检验了这个预言。
\index{几何工程理论|hyperpage}\index{Seiberg--Witten曲线@Seiberg--Witten 曲线|hyperpage}
